\chapter{Character Advancement and Capabilities}
łabel{ch:advancement}

\begin{tcolorbox}[colback=blue!5!white, colframe=blue!75!black, title=Philosophy of Growth, fonttitle=\bfseries]
∈dex{advancement!philosophy}
\textbf{Every XP spent is a choice about who your character becomes −− invest in yourself, your network, or your legacy.}

Your character is not just a collection of numbers; they are a living story. When you spend XP on an Attribute, you are saying ``I become stronger, smarter, more resilient.'' When you invest in a Follower or an Asset, you are saying ``I build relationships and influence that outlast me.'' When you buy a Talent, you are saying ``This is my signature, my unique edge.''

\medskip
\noindent There is no single ``correct'' way to grow. A lone duelist (see \ref{sec:player−archetypes}) might pour everything into personal skill; a merchant prince might build a trading empire. Both are valid. The question is: \textit{what kind of story do you want to tell?}
\end{tcolorbox}

\section*{Why This Chapter Matters}
∈dex{advancement}∈dex{Experience Points (XP)}

In this game, growth is not just about numbers −−− it is about deciding who your
character becomes. Advancement through \textbf{Experience Points (XP)} lets you
shape your capabilities, influence, and legacy in the world. Every XP choice is
a statement about your character's priorities and the mark they leave behind.

Your capabilities are built on four core \textbf{Attributes} and a set of
specialized \textbf{Skills}. This chapter explains how they work together,
what advancement looks like in play, and offers player‑facing tips for making
satisfying choices.

% ============================================================
% CORE ATTRIBUTES
% ============================================================
\section{Core Attributes}
łabel{sec:core−attributes}
∈dex{Attributes}∈dex{Body}∈dex{Wits}∈dex{Spirit}∈dex{Presence}

Attributes represent your character's fundamental capabilities. Each is rated
from 1 to 5; higher numbers mean greater potential in that area.

\subsection{Body}
łabel{sec:attribute−body}
∈dex{Body}
Physical strength, endurance, coordination, and health.
\begin{itemize}
  ℹtem \textbf{Used for:} melee combat, athletics, endurance tests, physical labor.
  ℹtem \textbf{Typical applications:} lifting, running, climbing, fighting, resisting physical harm.
  ℹtem \textbf{Associated skills:} Athletics, Brawl, Melee, Endurance (see \ref{sec:skills}).
\end{itemize}
\textbf{Rating guide:}
\begin{itemize}
  ℹtem 1: Average person; some physical activity.
  ℹtem 2: Fit; trains or works physically on a regular basis.
  ℹtem 3: Athlete or soldier; excellent condition.
  ℹtem 4: Exceptional athlete; near peak human.
  ℹtem 5: Peak human capability; legendary strength and stamina.
\end{itemize}

\subsection{Wits}
łabel{sec:attribute−wits}
∈dex{Wits}
Mental acuity, perception, quick thinking, and problem−solving.
\begin{itemize}
  ℹtem \textbf{Used for:} investigation, perception, tactics, snap decisions.
  ℹtem \textbf{Typical applications:} spotting details, solving puzzles, planning, reacting quickly.
  ℹtem \textbf{Associated skills:} Investigation, Tactics, Lore, Survival (see \ref{sec:skills}).
\end{itemize}
\textbf{Rating guide:}
\begin{itemize}
  ℹtem 1: Average awareness; sometimes misses important cues.
  ℹtem 2: Observant; notices most relevant details.
  ℹtem 3: Sharp−minded; quick to spot patterns.
  ℹtem 4: Exceptionally perceptive; rarely surprised.
  ℹtem 5: Almost uncanny awareness; sees connections others miss.
\end{itemize}

\subsection{Spirit}
łabel{sec:attribute−spirit}
∈dex{Spirit}
Willpower, intuition, mental resilience, and connection to intangible forces.
\begin{itemize}
  ℹtem \textbf{Used for:} resisting mental effects, intuition, magical aptitude, persistence.
  ℹtem \textbf{Typical applications:} resisting fear, sensing danger, channeling magic, enduring hardship.
  ℹtem \textbf{Associated skills:} Arcana, Ritual (for magic users), Endurance, Insight.
\end{itemize}
\textbf{Rating guide:}
\begin{itemize}
  ℹtem 1: Average willpower; somewhat suggestible.
  ℹtem 2: Strong−minded; resists ordinary pressure.
  ℹtem 3: Very determined; hard to intimidate.
  ℹtem 4: Exceptional will; can inspire others.
  ℹtem 5: Iron will; nearly unshakeable resolve.
\end{itemize}

\subsection{Presence}
łabel{sec:attribute−presence}
∈dex{Presence}
Charisma, social influence, appearance, and force of personality.
\begin{itemize}
  ℹtem \textbf{Used for:} social interactions, leadership, persuasion, intimidation.
  ℹtem \textbf{Typical applications:} negotiating, leading, charming, commanding attention.
  ℹtem \textbf{Associated skills:} Sway, Command, Performance, Deception (see \ref{sec:skills}).
\end{itemize}
\textbf{Rating guide:}
\begin{itemize}
  ℹtem 1: Average presence; does not particularly stand out.
  ℹtem 2: Noticeable; leaves a mild impression.
  ℹtem 3: Charismatic; naturally influential.
  ℹtem 4: Commanding; people listen when you speak.
  ℹtem 5: Magnetic; can sway crowds and redefine the mood of a room.
\end{itemize}

% ============================================================
% SKILL SYSTEM (REVISED)
% ============================================================
≠wpage
\section{Skill System}
łabel{sec:skills}
∈dex{Skills}

Skills represent specialized training and expertise. In play, you combine an
Attribute with a relevant Skill to form your dice pool (see \ref{ch:core−mechanics}).

\subsection{Skill Ratings}
łabel{sec:skill−ratings}
∈dex{Skill levels}
\begin{center}
\begin{tabular}{cl}
⊤rule
\textbf{Rating} & \textbf{Description} \\
∣rule
0 & Untrained −−− no formal experience \\
1 & Novice −−− basic understanding \\
2 & Competent −−− reliable and capable \\
3 & Professional −−− expert in the field \\
4 & Master −−− renowned specialist \\
5 & Grand Master −−− legendary skill \\
⊥tomrule
\end{tabular}
\end{center}

\subsection{Skill Categories}
łabel{sec:skill−categories}
∈dex{Combat skills}∈dex{Physical skills}∈dex{Social skills}∈dex{Knowledge skills}∈dex{Magic skills}
Below are the eighteen core skills, grouped by theme. Your table may rename or reskin
these, but the roles remain similar.

\paragraph{Combat Skills}
\begin{itemize}
  ℹtem \textbf{Melee}: swords, axes, close−quarters weapons.
  ℹtem \textbf{Ranged}: bows, crossbows, thrown weapons, firearms (if present).
  ℹtem \textbf{Brawl}: unarmed combat, grappling, improvised fighting.
  ℹtem \textbf{Tactics}: battlefield strategy, unit coordination, ambush planning.
\end{itemize}

\paragraph{Physical Skills}
\begin{itemize}
  ℹtem \textbf{Athletics}: running, climbing, jumping, swimming.
  ℹtem \textbf{Stealth}: moving unseen, hiding, blending into shadows.
  ℹtem \textbf{Endurance}: resisting fatigue, harsh weather, poison, disease.
  ℹtem \textbf{Craft}: building, repairing, creating tools, art, or structures.
  ℹtem \textbf{Survival}: tracking, foraging, predicting weather, navigating wilderness, animal handling (Nature folded into this skill).
\end{itemize}

\paragraph{Social Skills}
\begin{itemize}
  ℹtem \textbf{Sway}: persuasion, negotiation, charm.
  ℹtem \textbf{Command}: leadership, intimidation, issuing orders.
  ℹtem \textbf{Deception}: lying, bluffing, misdirection.
  ℹtem \textbf{Performance}: entertainment, oration, acting, musical display (used by Cantors for Songs).
  ℹtem \textbf{Insight}: reading emotions, detecting lies, understanding motives (covers social observation and ``reading'' people).
\end{itemize}

\paragraph{Knowledge Skills}
\begin{itemize}
  ℹtem \textbf{Lore}: history, culture, religion, customs, general knowledge (cities, civilizations, scholarly topics).
  ℹtem \textbf{Investigation}: research, deduction, analysis of clues (covers finding and interpreting evidence).
  ℹtem \textbf{Medicine}: healing, anatomy, first aid, treatment.
  % Nature is folded into Survival; no separate Nature skill.
\end{itemize}

\paragraph{Magic Skills}
\begin{itemize}
  ℹtem \textbf{Arcana}: magic theory, rituals (identification), mystical phenomena, free casting (TAGS system), counter−spelling.
  ℹtem \textbf{Ritual}: performing rites, invocations, exorcisms, consecrations −− used by Runekeepers, Invokers, Cantors, and Summoners for active magical workings.
\end{itemize}

\noindent\textbf{Optional Module −− Psionics:} If your table uses the \textit{Psionics} expansion (see \textit{The Grey Wanderer's Grimoire}), a separate \textbf{Psionics} skill is added. It is not part of the core skill list and is used for telepathy, telekinesis, clairvoyance, etc.

% ============================================================
% BUILDING DICE POOLS (unchanged except skill names)
% ============================================================
≠wpage
\section{Building Dice Pools}
łabel{sec:dice−pools}
∈dex{Dice pools}

Your dice pool for any action is:
\[
\text{Dice Pool} = \text{Attribute} + \text{Skill}.
\]

You roll that many d10s and interpret the results according to the core rules
(see \ref{ch:core−mechanics}).

\subsection{Choosing the Right Combination}
łabel{sec:attribute−skill−pairs}
∈dex{Attribute!combinations}∈dex{Skill!combinations}

The same fictional action can often be approached with different Attribute + Skill
pairs, depending on how your character does it. The GM has final say, but the
default is to use the most obvious pairing (e.g., \textit{Body + Melee} for a
straightforward sword swing). Creative alternatives are encouraged when they
fit the fiction.

\begin{itemize}
  ℹtem \textbf{Climbing a wall:}
    \begin{itemize}
      ℹtem \textit{Body + Athletics}: sheer strength and stamina.
      ℹtem \textit{Wits + Athletics}: clever route‑finding and leverage.
      ℹtem \textit{Spirit + Athletics}: grim determination against fear or pain.
    \end{itemize}
  ℹtem \textbf{Persuading a guard:}
    \begin{itemize}
      ℹtem \textit{Presence + Sway}: charm and personality.
      ℹtem \textit{Wits + Sway}: logical, well‑constructed arguments.
      ℹtem \textit{Spirit + Sway}: conviction that burns through doubt.
    \end{itemize}
  ℹtem \textbf{Investigating a crime scene:}
    \begin{itemize}
      ℹtem \textit{Wits + Investigation}: methodical observation and deduction.
      ℹtem \textit{Spirit + Investigation}: intuitive leaps and gut feelings.
      ℹtem \textit{Presence + Investigation}: getting witnesses to open up.
    \end{itemize}
\end{itemize}

\subsection{Creative Combinations}
łabel{sec:creative−dice−pools}
With your GM's agreement, you can justify unusual combinations that fit the
fiction. The table below offers common alternatives, but the GM may allow
others based on the scene.

\begin{center}
\small
\begin{tabular}{lll}
⊤rule
\textbf{Intended Action} & \textbf{Standard Pair} & \textbf{Alternative Pairs (Example)} \\
∣rule
Recall a combat drill & Wits + Lore & Body + Lore (muscle memory) \\
Soothe a patient & Presence + Medicine & Spirit + Medicine (laying on hands) \\
Forge a masterwork & Wits + Craft & Spirit + Craft (inspired artistry) \\
Read weather signs & Wits + Survival & Spirit + Survival (omen‑sense) \\
Feint in a duel & Wits + Melee & Presence + Melee (intimidation feint) \\
Track a spirit & Wits + Survival & Spirit + Lore (follow its legend) \\
Bribe a guard & Presence + Sway & Wits + Streetwise (know the right price) \\
⊥tomrule
\end{tabular}
\end{center}

\subsection{GM Guidance for Creative Pairs}
łabel{sec:gm−creative−pairs}
\begin{itemize}
  ℹtem \textbf{Ask ``How?''} −− The player must describe \textit{how} they use
        that Attribute and Skill together.
  ℹtem \textbf{Maintain Position} −− A creative pair does not automatically
        improve Position or reduce DV. It just changes the dice pool.
  ℹtem \textbf{Consistency} −− If you allow \textit{Spirit + Craft} for
        visionary work, allow it again in similar circumstances.
  ℹtem \textbf{No Double Dipping} −− The player cannot swap both Attribute
        \textit{and} Skill to max out a single stat; the combination must
        still respect the fiction.
\end{itemize}

\subsection{When to Roll (and When Not To)}
łabel{sec:when−to−roll}
Not every action requires a roll. The GM should call for a roll when:
\begin{enumerate}
  ℹtem The outcome is uncertain.
  ℹtem Failure would meaningfully change the situation.
  ℹtem The action is significant (not routine).
\end{enumerate}
For routine actions with no real risk, the group can simply agree that you
succeed. Example: climbing a knotted rope in calm weather −− no roll.
Climbing the same rope while goblins shoot arrows −− roll.

\subsection{Assistance and Teamwork}
łabel{sec:assistance−teamwork}
When multiple characters contribute to the same action:
\begin{itemize}
  ℹtem \textbf{Assistance:} One character is primary; up to three others add
        +1 die each by describing how they help (e.g., ``I hold the ladder steady'').
  ℹtem \textbf{Cooperation:} Several characters each roll; the group succeeds
        or fails based on combined results (e.g., all must sneak past a patrol).
  ℹtem \textbf{Complementary:} Different skills cover different parts of a
        larger task (e.g., one distracts a guard, another picks the lock).
\end{itemize}

\emph{Example:} A ranger scales an ice wall using \textbf{Wits + Athletics}
to find the safest holds, then switches to \textbf{Body + Athletics} to muscle over the
lip. The fiction guides which combination fits each moment.

\subsection{Selling the Action −− Intricate Description}
łabel{sec:intricate−description}
∈dex{Intricate description}∈dex{Selling it}

You are always welcome to state your action plainly. But when you take a moment
to \textbf{sell it} −− to invest in the fiction, reveal what your character cares
about, or show why this moment matters −− the GM may reward you with a mechanical
edge. This is about \emph{roleplay investment}, not word count.

\begin{itemize}
  ℹtem \textbf{Basic description:} ``I attack the goblin with my sword.''
        No bonus.

  ℹtem \textbf{Tactical description:} Use the environment, positioning, or a
        clever tactic. ``I feint high, then drive my shoulder into the goblin's
        shield, shoving it aside to open its flank.''
        The GM may or may not upgrade your Position by one step (e.g., \textit{Controlled}
        $→$ \textit{Dominant}) for this roll.

  ℹtem \textbf{Character‑driven description:} Reveal a stake, a bond, a fear,
        or a hard choice. \emph{One heartfelt sentence counts more than a
        paragraph of adjectives.}
        \begin{itemize}
          ℹtem Example: ``I think of my mother, who died because I wasn't brave
                enough to fight. I won't let that happen again. I plant my feet
                and dare the goblin to come through me.''
        \end{itemize}
        The GM may grant a larger benefit −− e.g., upgrade Position \emph{and}
        bank a \textbf{Boon} you can use later, or upgrade Effect by one step.
\end{itemize}

⫕section*{What Counts (and What Doesn't)}
\textbf{Counts:}
\begin{itemize}
  ℹtem Stating a clear stake: ``If I fail this lockpick, my sister doesn't eat.''
  ℹtem Invoking a bond mid‑action: ``I owe Ethan my life. I'm not letting him
        face this alone.''
  ℹtem Revealing a fear or hope tied to the roll: ``I'm terrified of fire, but
        I grab the burning beam −− because the children are behind me.''
  ℹtem Making a hard choice visible: ``I could run, but I don't. Running is what
        got me exiled.''
\end{itemize}
\textbf{Does not count:}
\begin{itemize}
  ℹtem Piling on adjectives (``gleaming,'' ``mighty,'' ``ancient'') without stakes.
  ℹtem Passive sensory laundry lists (``the cold wind, the wet stones, the dim light'').
  ℹtem Repeating the same florid description for every roll.
\end{itemize}

⫕section*{Guidelines for the GM}
\begin{itemize}
  ℹtem \textbf{No penalty for plain speech.} Never punish a player who is tired,
        shy, or simply prefers quick resolution. The bonus is a gift, not a tax.
  ℹtem \textbf{Be consistent.} Reward genuine investment, not vocabulary. A
        quiet player who delivers one heartfelt sentence deserves the bonus as
        much as an eloquent storyteller.
  ℹtem \textbf{The bonus should fit the fiction.} A desperate, reckless action
        described with real stakes might upgrade Effect but worsen Position −−
        let the mechanics mirror the emotion.
  ℹtem \textbf{Invite quiet players.} Ask leading questions: ``What's at stake
        for your character right now?'' or ``What are you afraid will happen if
        you fail?'' Gentle prompts help shy players find their moment.
  ℹtem \textbf{Group ``selling.''} If the whole table contributes to describing a
        group action (a coordinated ambush, a joint ritual), grant the bonus to
        the primary roller or give each participant a small benefit.
\end{itemize}

⫕section*{Example}
\textit{Basic:} ``I sneak past the sentry.'' (No bonus, roll as normal.)

\textit{Selling it (tactical):} ``I wait for the wind to rattle the shutters, then
pad across the wet cobblestones −− heel‑toe, heel‑toe −− keeping the stack of
barrels between me and the sentry's gaze.'' The GM nods: ``Great. Because you're
using the wind and the barrels, roll at \textit{Dominant} Position.''

\textit{Selling it (character‑driven):} ``My little brother is in that camp.
I've hidden from every fight since we left home −− but not this time. I creep
forward, not quiet anymore, because he needs to know someone came for him.''
The GM: ``That's the heart of it. Roll at \textit{Dominant} Position, and I'll
bank a Story Beat to make his reaction matter.''

\noindent
\textit{Remember:} The dice still decide success or failure. But a well‑told
action makes the story richer −− and earns you an edge when the dice fall.

% ============================================================
% EARNING XP (unchanged)
% ============================================================
≠wpage
\section{Earning Experience Points}
łabel{sec:earning−xp}
∈dex{Experience Points (XP)!earning}

XP represents learning through action. You gain it by engaging meaningfully
with the world: overcoming challenges, taking risks, making hard choices, and
leaning into your character's story.

\subsection{Session Breakdown}
łabel{sec:session−xp}
At the end of each session, XP is awarded based on what happened at the table.
A typical breakdown might look like:

\begin{itemize}
  ℹtem \textbf{Base Participation}: +2~XP for attending and contributing.
  ℹtem \textbf{Major Objectives}: +2−−4~XP for completing significant story goals.
  ℹtem \textbf{Discoveries}: +1−−2~XP for uncovering important lore, locations, or secrets.
  ℹtem \textbf{Hard Choices}: +1−−2~XP for making difficult moral or strategic decisions.
  ℹtem \textbf{Story Engagement}: +1−−3~XP for embracing complications and twists.
  ℹtem \textbf{Personal Goals}: +1−−2~XP for pursuing your character's individual storylines.
\end{itemize}

\emph{Example:} At the end of a session, the party rogue gains:
\begin{itemize}
  ℹtem +2 XP for participation,
  ℹtem +2 XP for helping retrieve a cursed artifact,
  ℹtem +1 XP for pushing a personal rivalry subplot.
\end{itemize}
Total: 5 XP.

\subsection{Game Pace Options}
łabel{sec:game−pace}
∈dex{Game pace}
You and your GM can decide how quickly characters advance, matching the tone
you want:

\begin{center}
\begin{tabular}{cll}
⊤rule
\textbf{Mode} & \textbf{XP / Session} & \textbf{Tone} \\
∣rule
Gritty   & 4−−6  & Hard choices, slow growth \\
Standard & 6−−10 & Balanced progression \\
Epic     & 10−−14 & Heroic, rapid development \\
⊥tomrule
\end{tabular}
\end{center}

\noindent\textbf{Player Tip:} If you want a sweeping, mythic tale, ask for an
\emph{Epic} pace. If you want advancement to feel hard‑won, lean toward
\emph{Gritty}.

\subsection{Arc Completion Bonus}
łabel{sec:arc−bonus}
When a major story arc finishes (typically 3−−6 sessions), everyone gains an
additional +8−−12 XP. One player may receive +2 XP more for a particularly
memorable contribution. This rewards the story as a whole, not just individual
rolls.

\begin{tcolorbox}[colback=yellow!5!white, colframe=yellow!75!black, title=Boons vs. XP]
∈dex{Boons!conversion}
You may convert 2 Boons $→$ 1 XP (once per session, max 2 XP per session via conversion). \textbf{But:} Spending Boons during play is almost always more fun. A re‑rolled failure, an improved Position, or an activated Asset creates a memorable moment. Hoarding Boons for XP conversion starves the table of drama. Spend freely; the XP will come from your adventures.
\end{tcolorbox}

% ============================================================
% SPENDING XP (unchanged except skill names)
% ============================================================
\section{Spending Experience Points}
łabel{sec:spending−xp}
∈dex{Experience Points (XP)!spending}

XP is your currency for growth. You can invest it in three broad areas:

\begin{enumerate}
  ℹtem \textbf{Personal Improvement} (Attributes and Skills)
  ℹtem \textbf{Special Abilities} (Talents, magic, unique tricks)
  ℹtem \textbf{Resources and Influence} (Assets, followers, organizations)
\end{enumerate}

Think of these as three pillars: \textit{what you can do}, \textit{what makes you unique},
and \textit{what you own or command}.

\subsection{Personal Improvement}

⫕section*{Attributes}
∈dex{Attributes!increasing}
Attributes are expensive but powerful. The cost to raise an Attribute is:
\[
\text{XP Cost} = 3 × \text{new rating}.
\]
Examples:
\begin{itemize}
  ℹtem Raising Body from 2 to 3 costs $3 × 3 = 9$ XP.
  ℹtem Raising Spirit from 4 to 5 costs $3 × 5 = 15$ XP.
\end{itemize}
Attribute increases usually require downtime equal to the new rating in days,
spent training, reflecting, or transforming yourself (see \ref{sec:training−time}).

⫕section*{Skills}
∈dex{Skills!increasing}
Skills are cheaper and represent focused practice. The cost is:
\[
\text{XP Cost} = 2 × \text{new level}.
\]
Examples:
\begin{itemize}
  ℹtem Lore 1 $→$ 2 costs $2 × 2 = 4$ XP.
  ℹtem Melee 3 $→$ 4 costs $2 × 4 = 8$ XP.
\end{itemize}
Skill increases typically require downtime equal to the new level in days.

\medskip
\noindent\textit{Example:} Kara wants to improve her swordsmanship (Melee) from 2 to 3.  
She saves 6 XP and spends three in‑game days training with her mentor. This
training becomes part of the story, not just a line on a sheet.

\subsection{Special Abilities}
∈dex{Talents}∈dex{Special abilities}
Special abilities are unique moves, talents, or tricks that define your style.

\begin{description}
  ℹtem[General Abilities (4−−8 XP)]  
    Universal benefits, like improved recovery, bonus dice in a common
    situation, or a signature combat technique.  
    \emph{Example:} ``Quick Recovery'' −−− heal 1 extra Harm when you rest.

  ℹtem[Cultural Abilities (variable)]  
    Abilities tied to heritage or background. They may require specific
    fictional positioning (see \ref{sec:small−folk−detailed}).  
    \emph{Example:} ``Stone Sense'' −−− an instinctive feel for stonework and tunnels.

  ℹtem[Advanced Abilities (12+ XP)]  
    High‑impact features available at higher tiers, often with strong narrative
    hooks or prerequisites (see \ref{sec:tiers}).  
    \emph{Example:} ``Master Diplomat'' −−− once per session, you can reroll or
    soften the outcome of a failed major negotiation.
\end{description}

\emph{Example:} A veteran bard purchases \emph{Silver Tongue} for 6 XP,
gaining the ability to sway a hostile crowd once per session. This quickly
becomes their trademark in tense public scenes.

\subsection{Resources and Influence}
∈dex{Assets}∈dex{Followers}
See §\ref{ch:assets−followers} for full rules. As a quick reference:

\begin{description}
  ℹtem[Minor Resource (4 XP, ~1 week)]  
    A small shop, basic workshop, discreet safe house, or minor contact
    network. Offers modest but reliable benefits.

  ℹtem[Standard Resource (8 XP, ~2 weeks)]  
    A decent‑sized business, crew, or operation with real impact. Comes with
    some upkeep and obligations.

  ℹtem[Major Resource (12 XP, ~1 month)]  
    A large enterprise, elite team, or rare capability. Very potent but also
    demanding; often becomes a major part of your character's story.
\end{description}

\noindent\textbf{Player Tip:} Resources pull the camera out from your character
to your \emph{sphere of influence}. A spy network means intrigue. A workshop
means invention. A guild hall means politics.

\begin{tcolorbox}[colback=green!5!white, colframe=green!75!black, title=When to Invest in Assets and Followers]
∈dex{Assets!investment advice}∈dex{Followers!investment advice}
\begin{center}
\begin{tabular}{lll}
⊤rule
\textbf{Player Archetype} & \textbf{Typical Investment} & \textbf{Best For} \\
∣rule
Solo & 0−−10% XP & Reliable personal power, minimal upkeep. \\
Mixed (Balanced) & 15−−░ XP & One solid Follower or Asset (e.g., a scout or a safehouse). \\
Mastermind & 35−−55% XP & Networks, informants, multiple Assets; build a faction. \\
⊥tomrule
\end{tabular}
\end{center}
\textbf{Planning for succession?} A loyal Follower (Cap 3+) can become your successor −− see the \textbf{Legacy Engine} (§\ref{ch:legacy−engine}). Investing in a potential heir early pays off when your character retires or falls.
\end{tcolorbox}

% ============================================================
% DEVELOPMENT PATHS (unchanged)
% ============================================================
≠wpage
\section{Character Development Paths}
łabel{sec:development−paths}
∈dex{Development paths}

How you spend XP shapes your character's growth. These are broad patterns, not
rigid classes.

\begin{center}
\begin{tabular}{p{3cm}p{2.8cm}p{6cm}}
⊤rule
\textbf{Path} & \textbf{Typical XP Split} & \textbf{Strengths / Weaknesses} \\
∣rule
Specialist & 70−−90% personal, 0−−10% resources, 0−−20% abilities & Extremely effective in a narrow niche. Vulnerable outside it. \\
Leader     & 50−−65% personal, 15−−░ resources, 15−−░ abilities & Well‑rounded, boosts the group. Rarely the best at one thing. \\
Mastermind & 25−−40% personal, 35−−55% resources, 20−−40% abilities & Solves problems through assets. Personally vulnerable, many moving parts. \\
⊥tomrule
\end{tabular}
\end{center}

\noindent\textit{Example (Leader):} A merchant‑prince spends 60% of XP on personal
skills and social abilities, 20% on a network of trade contacts, and 20% on
talents that improve negotiation and command.

These patterns are \emph{guides}, not rules. You can shift between them as your
story evolves.

% ============================================================
% TRAINING & TIME (unchanged)
% ============================================================
\section{Training and Development Time}
łabel{sec:training−time}
∈dex{Training}∈dex{Downtime}

Most improvements take time in the fiction. Training montages, study scenes,
and building projects are all opportunities for character moments.

\subsection{Standard Time Requirements}
łabel{sec:standard−training}
\begin{itemize}
  ℹtem \textbf{Attribute increase:} new rating in days.
  ℹtem \textbf{Skill improvement:} new level in days.
  ℹtem \textbf{Resource acquisition:} about 1 week to 1 month depending on scope.
  ℹtem \textbf{Ability learning:} typically 3−−10 days.
\end{itemize}

\subsection{Accelerated Development (Rush Rule)}
łabel{sec:rush−rule}
∈dex{Rush rule}
You \emph{can} attempt to learn faster, but rushing has a price.
\begin{itemize}
  ℹtem The group tracks a short \textbf{Risk Timer } (4 segments) for your crash training.
  ℹtem If it fills, the new capability has flaws:
    \begin{itemize}
      ℹtem Attribute/Skill: temporary −−1 die penalty until you retrain properly.
      ℹtem Resource: loyalty issues, partial effectiveness, or hidden problems.
      ℹtem Ability: unreliable effects or quirky side consequences.
    \end{itemize}
\end{itemize}
\noindent\textit{Example:} A wizard crams advanced spellwork into three frantic days.
She gains the ability, but the Risk Timer fills: her spells spark and flicker
until she takes time to train safely.

% ============================================================
% TIERS (unchanged)
% ============================================================
≠wpage
\section{Character Progression Tiers}
łabel{sec:tiers}
∈dex{Tiers}∈dex{Progression}

As you accumulate XP and capabilities, you advance through tiers that represent your growing reputation and influence.

\subsection{Tier I: Novice (0−−40 XP)}
łabel{sec:tier−i}
\begin{itemize}
ℹtem Learning the ropes, establishing yourself.
ℹtem Local reputation, modest capabilities.
ℹtem Typical assets: basic equipment, a few contacts.
ℹtem \textbf{Position Manifestation:} The difference between life and death; a shift from Desperate to Controlled means ``I live to fight another day.''
\end{itemize}

\subsection{Tier II: Seasoned (41−−90 XP)}
łabel{sec:tier−ii}
\begin{itemize}
ℹtem Proven capability, reliable in your specialty.
ℹtem Regional reputation, respected in your field.
ℹtem Typical assets: skilled followers, specialized equipment.
ℹtem \textbf{Position Manifestation:} The difference between victory and defeat; a shift from Controlled to Dominant means ``I've earned the right to choose how this ends.''
\end{itemize}

\subsection{Tier III: Veteran (91−−150 XP)}
łabel{sec:tier−iii}
\begin{itemize}
ℹtem Master of your craft, significant influence.
ℹtem National reputation, can handle major challenges.
ℹtem Typical assets: multiple operations, elite teams.
ℹtem \textbf{Position Manifestation:} The difference between influence and control; a shift from Controlled to Dominant means ``The city hangs on my decision.''
\end{itemize}

\subsection{Tier IV: Paragon (151−−220 XP)}
łabel{sec:tier−iv}
\begin{itemize}
ℹtem Exceptional capability, strategic influence.
ℹtem Continental influence, shapes regional outcomes.
ℹtem Typical assets: organizations, unique capabilities.
ℹtem \textbf{Position Manifestation:} The difference between shaping events and defining eras; a shift from Controlled to Dominant means ``I am the axis upon which history turns.''
\end{itemize}

\subsection{Tier V: Mythic (221+ XP)}
łabel{sec:tier−v}
\begin{itemize}
ℹtem Legendary status, world‑changing influence.
ℹtem Historical reputation, defines eras.
ℹtem Typical assets: nations, legendary artifacts.
ℹtem \textbf{Position Manifestation:} The difference between ending and remaking the world; a shift from Desperate to Controlled means ``The universe waits on my choice.''
\end{itemize}

\subsection{Position Scales with Narrative Weight}
łabel{sec:position−narrative−weight}
The mechanical effect of Position (re‑rolling one die, see \ref{ch:core−mechanics}) never changes −−− but what it \textit{means} evolves with your tier:
\begin{itemize}
ℹtem \textbf{Tier I:} Position defines personal survival.
ℹtem \textbf{Tier III:} Position shapes regional outcomes.
ℹtem \textbf{Tier V:} Position reshapes reality.
\end{itemize}
The system does not need to scale −−− your story does.

\begin{tcolorbox}[colback=blue!5!white, colframe=blue!75!black, boxrule=0.5pt, arc=2mm]
\textbf{GM Tip:} Never add mechanical bonuses to Position. Instead, deepen the narrative weight:
\begin{itemize}
ℹtem \textbf{Tier I:} ``Your Position shows whether you live or die.''
ℹtem \textbf{Tier V:} ``Your Position shows how the world bends to your will.''
\end{itemize}
The mechanics stay consistent so the narrative can grow as large as your imagination allows.
\end{tcolorbox}

% ============================================================
% ALLIES & FOLLOWERS (unchanged)
% ============================================================
≠wpage
\section{Managing Allies and Followers}
łabel{sec:allies−followers}
∈dex{Allies}∈dex{Followers}

Allies can extend your reach but also create obligations and risks. For full
rules, see §\ref{ch:assets−followers}.

\subsection{Acquisition Costs}
łabel{sec:follower−cost}
A skilled helper's cost in XP is roughly their \emph{capability rating squared}.
\emph{Example:} A capability 3 scout costs $3² = 9$ XP.

\subsection{Upkeep Requirements}
łabel{sec:follower−upkeep}
Each downtime period, you either:
\begin{itemize}
  ℹtem spend XP equal to their capability rating, \textbf{or}
  ℹtem dedicate a meaningful scene to maintaining the relationship (time, favors, protection).
\end{itemize}

\subsection{Risk Management}
łabel{sec:follower−risk}
\begin{itemize}
  ℹtem When big complications arise, your allies can be targeted instead of you.
  ℹtem Allies can handle one significant task off‑screen per downtime, but this often creates extra story fallout (the GM gains a Story Beat).
\end{itemize}

% ============================================================
% SKILL ADVANCEMENT & SYNERGY (updated)
% ============================================================
\section{Skill Advancement and Synergies}
łabel{sec:skill−advancement}
∈dex{Skills!advancement}∈dex{Synergy}

\subsection{XP Costs}
\begin{center}
\begin{tabular}{cl}
⊤rule
\textbf{Improvement} & \textbf{XP Cost} \\
∣rule
0 $→$ 1 & 2 XP \\
1 $→$ 2 & 4 XP \\
2 $→$ 3 & 6 XP \\
3 $→$ 4 & 8 XP \\
4 $→$ 5 & 10 XP \\
⊥tomrule
\end{tabular}
\end{center}

\subsection{Training Time}
łabel{sec:skill−training−time}
\begin{itemize}
  ℹtem 0 $→$ 1: 1 day of practice.
  ℹtem 1 $→$ 2: 3 days of training.
  ℹtem 2 $→$ 3: 1 week of intensive study.
  ℹtem 3 $→$ 4: 2 weeks with a master or rigorous practice.
  ℹtem 4 $→$ 5: 1 month of dedicated focus.
\end{itemize}

\subsection{Attribute Limits}
łabel{sec:attribute−limits}
You cannot have a Skill rating higher than the relevant Attribute. To advance
a Skill further, you must first raise its primary Attribute (see \ref{sec:core−attributes}).

\subsection{Skill Synergies}
łabel{sec:skill−synergies}
When the fiction supports it, well‑justified combinations of skills can grant
+1 die or better Position/Effect. Examples:
\begin{itemize}
  ℹtem \textbf{Combat:} Tactics + Command (battlefield coordination), Melee + Athletics (acrobatic strikes).
  ℹtem \textbf{Social:} Sway + Lore (persuasion backed by knowledge), Deception + Performance (long cons).
  ℹtem \textbf{Exploration:} Investigation + Survival (tracking), Lore + Insight (cultural observation).
\end{itemize}

% ============================================================
% USING SKILLS IN PLAY (updated)
% ============================================================
\section{Using Skills in Play}
łabel{sec:using−skills}
∈dex{Skills!using in play}

\subsection{When to Roll}
łabel{sec:when−to−roll}
You roll when:
\begin{itemize}
  ℹtem the outcome is uncertain,
  ℹtem failure would matter to the story, and
  ℹtem the action is significant enough to deserve the spotlight.
\end{itemize}
For routine actions with no real risk, the group can often simply agree that
you succeed.

\subsection{Difficulty and Skill Level}
łabel{sec:difficulty−skill}
As your skill rises, routine tasks become trivial. A rough guideline:

\begin{center}
\begin{tabular}{cll}
⊤rule
\textbf{Skill Level} & \textbf{Routine Task} & \textbf{Challenging Task} \\
∣rule
0 & Needs a roll; modest DV & Needs a roll; higher DV \\
1 & Often needs a roll & Challenging but fair \\
2 & Usually automatic & Needs a roll; moderate DV \\
3 & Automatic & Needs a roll; interesting stakes \\
4+ & Automatic & Automatic for simple cases \\
⊥tomrule
\end{tabular}
\end{center}

\noindent\textbf{Reading this as a player:} by Skill 3, basic tasks in your area of
expertise should rarely require rolls. You are there to be impressive, not to
trip over routine jobs.

\subsection{Group Skill Use}
łabel{sec:group−skills}
When multiple characters contribute:
\begin{itemize}
  ℹtem \textbf{Assistance:} One character is primary; up to three others add +1 die each by describing how they help (see \ref{ch:core−mechanics}).
  ℹtem \textbf{Cooperation:} Several characters each roll; the group succeeds or fails based on combined results.
  ℹtem \textbf{Complementary:} Different skills cover different parts of a larger task (e.g. one distracts, one picks a lock).
\end{itemize}

% ============================================================
% SKILL CHALLENGES (unchanged)
% ============================================================
\section{Skill Challenges}
łabel{sec:skill−challenges}
∈dex{Skill challenges}∈dex{Timers}

Some goals are too big for a single roll.

\subsection{Extended Tests}
łabel{sec:extended−tests}
For long, involved tasks:
\begin{itemize}
  ℹtem The group sets a progress timer (4−−8 segments).
  ℹtem Each significant success fills one or more segments.
  ℹtem Complications can add segments or create new problems.
\end{itemize}

\subsection{Complex Challenges}
łabel{sec:complex−challenges}
For multi‑step or multi‑skill situations:
\begin{itemize}
  ℹtem Different skills tackle different phases of the challenge.
  ℹtem Success in one area creates opportunities; failure creates obstacles.
\end{itemize}
\noindent\textit{Example (Heist):} \textbf{Stealth} to infiltrate, \textbf{Craft} to bypass
locks, \textbf{Investigation} to find the vault, \textbf{Deception} to mislead guards.
Each roll advances a \emph{Heist Timer }; complications add heat, suspicion, or
security changes.

% ============================================================
% STRATEGIC ADVANCEMENT (unchanged)
% ============================================================
\section{Strategic Advancement Considerations}
łabel{sec:advancement−strategy}
∈dex{Advancement!strategies}

\subsection{Early Game (Tiers I−−II)}
łabel{sec:early−game}
\begin{itemize}
  ℹtem Focus on survival and a clear niche.
  ℹtem Raise a key Attribute to 3 early.
  ℹtem Bring core Skills to 2−−3.
  ℹtem Establish one or two modest resources or safe havens.
\end{itemize}

\subsection{Mid Game (Tier III)}
łabel{sec:mid−game}
\begin{itemize}
  ℹtem Develop signature moves or combinations.
  ℹtem Expand your resource base.
  ℹtem Consider your role in the group and lean into it.
\end{itemize}

\subsection{Late Game (Tiers IV−−V)}
łabel{sec:late−game}
\begin{itemize}
  ℹtem Choose advanced abilities that reshape the story around you.
  ℹtem Build organizations, movements, or legacies.
  ℹtem Decide how your character will leave their mark on the world.
\end{itemize}

\begin{tcolorbox}[colback=orange!5!white, colframe=orange!75!black, title=Common Advancement Mistakes]
\noindent\textbf{Avoid these pitfalls:}
\begin{itemize}
  ℹtem \textbf{Over‑specializing too early:} A Melee 5 character with Body 3 is fragile and incapable outside combat.
  ℹtem \textbf{Forgetting training time:} You cannot learn a new skill in the middle of a dungeon without justification.
  ℹtem \textbf{Neglecting resources:} A few well‑chosen assets can solve problems no amount of personal skill can.
  ℹtem \textbf{Hoarding XP:} Spend your XP! Characters who never improve become irrelevant to the escalating threats.
\end{itemize}
\end{tcolorbox}

% ============================================================
% PRACTICAL EXAMPLES (updated skills)
% ============================================================
\section{Practical Examples}
łabel{sec:practical−examples}
∈dex{Examples}

\subsection{Combat Example}
łabel{sec:combat−example}
A warrior with Body 4 and Melee 3 attacks:
\begin{itemize}
  ℹtem Dice pool: $4 + 3 = 7$d10.
  ℹtem Opponent's defenses and the fiction help set the difficulty (see \ref{ch:core−mechanics}).
  ℹtem Strong successes translate into higher Effect or better Position.
\end{itemize}

\subsection{Social Example}
łabel{sec:social−example}
A diplomat with Presence 3 and Sway 2 negotiates:
\begin{itemize}
  ℹtem Dice pool: $3 + 2 = 5$d10.
  ℹtem Position might be \emph{Controlled} if circumstances favor them, or worse if under suspicion.
  ℹtem Consequences for failure could be mistrust, delays, or new complications.
\end{itemize}

\subsection{Exploration Example}
łabel{sec:exploration−example}
A scout with Wits 3 and Survival 2 searches for tracks:
\begin{itemize}
  ℹtem Dice pool: $3 + 2 = 5$d10.
  ℹtem Easier tasks (fresh tracks) are more forgiving; older or obscured tracks are harder.
  ℹtem Partial successes might reveal some information but not all, or come with side effects (e.g., you find the tracks but disturb a nest of predators).
\end{itemize}

\noindent Your Attributes and Skills determine not just \emph{if} you succeed, but
\emph{how} you tackle problems. Lean into combinations that express your
character's style.

% ============================================================
% QUICK REFERENCE BOXES (updated)
% ============================================================
\begin{tcolorbox}[colback=green!5!white,colframe=green!75!black,title=XP Planning Guide,fonttitle=\bfseries]
\textbf{Early Tier Priorities:}
\begin{itemize}
  ℹtem Raise a core Attribute to 3 (9 XP).
  ℹtem Bring key Skills to 2−−3 (4−−6 XP each).
  ℹtem Claim 1−−2 minor resources (4 XP each).
\end{itemize}

\textbf{Mid Tier Expansion:}
\begin{itemize}
  ℹtem Push core Attributes toward 4 (12 XP).
  ℹtem Specialize key Skills to 4 (8 XP).
  ℹtem Add standard resources (8 XP each).
  ℹtem Explore cultural or thematic abilities (6−−10 XP).
\end{itemize}

\textbf{Late Tier Mastery:}
\begin{itemize}
  ℹtem Purchase capstone abilities (12+ XP).
  ℹtem Establish major resources (12 XP).
  ℹtem Launch legacy projects, organizations, or movements.
\end{itemize}
\end{tcolorbox}

\begin{tcolorbox}[colback=blue!5!white,colframe=blue!75!black,title=Attributes and Skills Quick Reference,fonttitle=\bfseries]
\textbf{Attributes (1−−5):}
\begin{itemize}
  ℹtem \textbf{Body}: physical capability
  ℹtem \textbf{Wits}: mental acuity
  ℹtem \textbf{Spirit}: willpower and inner strength
  ℹtem \textbf{Presence}: social influence
\end{itemize}

\textbf{Skill Levels:}
\begin{itemize}
  ℹtem 0: Untrained   1: Novice   2: Competent
  ℹtem 3: Professional   4: Master   5: Grand Master
\end{itemize}

\textbf{Core Skills (18):}
\begin{itemize}
  ℹtem Combat: Melee, Ranged, Brawl, Tactics
  ℹtem Physical: Athletics, Stealth, Endurance, Craft, Survival
  ℹtem Social: Sway, Command, Deception, Performance, Insight
  ℹtem Knowledge: Lore, Investigation, Medicine
  ℹtem Magic: Arcana, Ritual
\end{itemize}

\textbf{Dice Pool:} Attribute $+$ Skill d10s.

\textbf{Skill Improvement:} New level $×$ 2 XP (plus training time).

\textbf{Specialization:} At Skill 3+, you are a known expert in that field.

\textbf{Synergy:} Well‑justified combinations of skills can grant +1 die or
better Position/Effect when the fiction supports it.
\end{tcolorbox}